\documentclass[11pt]{article}
\usepackage[margin=1in]{geometry}
\usepackage[T1]{fontenc}
\usepackage{lmodern}
\usepackage{xcolor}
\usepackage{hyperref}
\usepackage{booktabs}
\usepackage{array}
\usepackage{enumitem}
\usepackage{titlesec}

\hypersetup{colorlinks=true,urlcolor=blue,linkcolor=blue}

\definecolor{reviewGreen}{RGB}{40,140,70}
\definecolor{oopBlue}{RGB}{40,90,180}
\definecolor{algoOrange}{RGB}{220,120,20}
\definecolor{dataPurple}{RGB}{120,50,160}

\title{\textbf{CPSC 250: Programming for Data Manipulation}\\
\large Course Roadmap and Zybooks Structure}
\author{Christopher Newport University}
\date{Summer 2026}

\begin{document}
\maketitle

\section{Course Philosophy}

CPSC 250 bridges introductory Python programming and modern computational problem solving. The course begins by reinforcing material from CPSC 150/150L, then moves into object-oriented programming, algorithms, recursion, scientific computing, data analysis, and visualization.

The goal is not simply to learn more Python syntax. The goal is to develop lasting computational problem-solving skills.

\section{How the CPSC 250 Zybook Is Organized}

The CPSC 250 Zybook is custom organized from material students may recognize from CPSC 150, but the ordering and emphasis have been changed to support the goals of this course.

Some early material is review and reinforcement. Later chapters introduce major new topics.

\subsection*{\color{reviewGreen}Review and Reinforcement from CPSC 150}

\begin{center}
\renewcommand{\arraystretch}{1.25}
\begin{tabular}{p{0.20\textwidth} p{0.35\textwidth} p{0.25\textwidth}}
\toprule
\textbf{CPSC 150 Zybook} & \textbf{Topic} & \textbf{CPSC 250 Placement} \\
\midrule
Ch. 1 & Introduction & Ch. 1 \\
Ch. 3 & Variables / Expressions & Ch. 2 \\
Ch. 4 & Primitive Types & Ch. 3 \\
Ch. 6 & Branching & Ch. 4 \\
Ch. 7 & Loops & Ch. 5 \\
Ch. 5 and Ch. 8 & Functions & Ch. 6 \\
Ch. 9 & Strings & Ch. 7 \\
Ch. 10 & Lists / Dictionaries & Ch. 8 \\
Ch. 12 & Modules & Ch. 11 \\
Ch. 11 & Files & Ch. 12 \\
\bottomrule
\end{tabular}
\end{center}

\noindent
This material is not intended to completely reteach CPSC 150. Instead, it reinforces the programming tools students are expected to carry forward.

\subsection*{\color{oopBlue}Major New CPSC 250 Topics}

\begin{center}
\renewcommand{\arraystretch}{1.25}
\begin{tabular}{p{0.25\textwidth} p{0.55\textwidth}}
\toprule
\textbf{CPSC 250 Chapter} & \textbf{Topic} \\
\midrule
Ch. 9 & Classes and Objects \\
Ch. 13 & Inheritance and Polymorphism \\
Ch. 14 & Recursion \\
Ch. 15 & Searching and Sorting Algorithms \\
Ch. 30 & Scientific Computing and Data Analysis \\
\bottomrule
\end{tabular}
\end{center}

\section{Course Structure}

\subsection*{\color{reviewGreen}Section I --- Review and Reinforcement}

\begin{itemize}
    \item variables, expressions, and primitive data types
    \item branching and loops
    \item functions
    \item strings, lists, and dictionaries
    \item files and modules
\end{itemize}

\subsection*{\color{oopBlue}Section II --- Object-Oriented Programming}

\begin{itemize}
    \item classes and objects
    \item constructors
    \item instance variables
    \item methods
    \item encapsulation
    \item operator overloading
\end{itemize}

\subsection*{\color{oopBlue}Section III --- Inheritance and Polymorphism}

\begin{itemize}
    \item inheritance
    \item derived classes
    \item method overriding
    \item polymorphism
\end{itemize}

\subsection*{\color{algoOrange}Section IV --- Algorithms and Problem Solving}

\begin{itemize}
    \item recursion
    \item searching algorithms
    \item sorting algorithms
    \item problem decomposition
    \item computational efficiency
\end{itemize}

\subsection*{\color{dataPurple}Section V --- Scientific Computing and Data Analysis}

\begin{itemize}
    \item NumPy
    \item pandas
    \item matplotlib
    \item data visualization
    \item curve fitting
    \item data analysis
\end{itemize}

\section{Lecture and Laboratory Relationship}

CPSC 250 and CPSC 250L are designed to complement one another.

\textbf{CPSC 250} focuses on concepts, computational thinking, algorithms, problem solving, and scientific/data-oriented programming.

\textbf{CPSC 250L} emphasizes active in-class programming, debugging, GitHub workflows, feature branches, version control, command-line literacy, and professional software development practices.

\section{Development Environment}

Students will use:

\begin{itemize}
    \item Python 3
    \item PyCharm Community Edition
    \item GitHub
    \item Git
\end{itemize}

Students enrolled in either CPSC 250 or CPSC 250L should complete the development environment setup process before the end of the first week.

\section{Summer Course Rhythm}

The summer course is highly compressed and moves quickly.

\begin{center}
\renewcommand{\arraystretch}{1.25}
\begin{tabular}{ll}
\toprule
Morning & CPSC 250 lecture \\
Afternoon & CPSC 250L laboratory \\
Fridays & Tests / assessments \\
June 19 & No classes, Juneteenth \\
\bottomrule
\end{tabular}
\end{center}

\section{Grading and Course Policies}

Detailed grading policies, assignment policies, attendance expectations, and course requirements are described in the official course syllabi.

\textbf{The syllabi are the authoritative source} for grading weights, assignment policies, late work policies, attendance expectations, academic integrity policies, and laboratory requirements.

\section{Final Thoughts}

Programming is learned by actively writing, testing, debugging, and refining code.

Students are encouraged to experiment with code examples, ask questions frequently, practice regularly, become comfortable debugging, and develop confidence with professional software tools.

\end{document}
